
Einleitung

Aktuelle Situation: primitive + zusammengesetzte Datentypen

Abb.: Puzzle vs. CYBOP-Schema

Alle Arten von Programmen möglich
OOP etabliert
Wozu über Alternativen nachdenken?

Auch funktionale Programmierung basiert vom Prinzip her
auf ähnlichen Konzepten, bietet lediglich eine andere Syntax.
(TODO, Nachweis, Beispiel)

Probleme: enge Kopplung (Verdrahtung) verursacht hohe Komplexität
(seit Jahrzehnten kritisiert).

Vielleicht lohnt es sich also doch, über Alternativen nachzudenken?

Wäre es trotz der Vielgestaltigkeit unterschiedlicher Anwendungen möglich,
ihre Datenstrukturen zu vereinheitlichen?

Eine uniforme Datenstruktur könnte nicht nur viel effizienter verwendet werden,
sondern auch und vor allem die Komplexität von Software verringern helfen.

Diese Arbeit möchte ein universell nutzbares Schema zur Beschreibung
von Daten vorstellen, wie sie in Standard-Anwendungssoftware zum
Einsatz kommen.
